% ===========================================================================
\section{同构定理}\label{sec:iso-theorems}

\subsection{第一同构定理（同态基本定理）}

\begin{theorem}[环的第一同构定理]
设 $f\colon R\to R'$ 为环同态，则 $f$ 可分解为满同态、同构、单同态的
复合：
\[
  R \xrightarrow{\ \pi\ } R/\ker f \xrightarrow{\ \bar f\ }
  \Img f \xrightarrow{\ \iota\ } R',
\]
其中 $\pi\colon a\mapsto a+\ker f$ 是典范投影（满同态），
$\iota\colon\Img f\hookrightarrow R'$ 是包含映射（单同态），
$\bar f\colon a+\ker f\mapsto f(a)$ 是诱导同构。特别地，
\[
  R/\ker f \;\cong\; \Img f.
\]
\end{theorem}

\begin{proof}[证明要点]
$\bar f$ 的良好定义性：若 $a' = a + i$（$i\in\ker f$），则
\[
  f(a') = f(a) + f(i) = f(a) + 0 = f(a).
\]
$\bar f$ 显然满；又 $\bar f(a+\ker f) = 0 \iff f(a) = 0 \iff a\in\ker f$，
故单，从而是同构。
\end{proof}

\subsection{理想的运算}

\begin{proposition}[理想之交仍为理想]
任意多个理想的交仍是理想（由定义逐条验证即得）。
\end{proposition}

\begin{definition}[生成理想与主理想（Principal Ideal）]
设 $S\subseteq R$，所有包含 $S$ 的理想之交是包含 $S$ 的\keyword{最小
理想}，称为 \keyword{$S$ 生成的理想}，记 $\langle S\rangle$。由单个
元素生成的理想 $\langle a\rangle$ 称为\keyword{主理想}。
\end{definition}

\begin{proposition}[生成理想的显式刻画]
\[
  \langle S\rangle = \left\{ \sum_{i=1}^{n} x_i\, s_i\, y_i
    \ \middle|\ n\in\N,\ x_i, y_i\in R,\ s_i\in S \right\}.
\]
特别地，$\langle a\rangle = \left\{\sum_i x_i\, a\, y_i\right\}$。
若 $R$ 交换（含幺），则
\[
  \langle a\rangle = aR = \{\,ar : r\in R\,\}.
\]
\end{proposition}

\begin{proof}
右端 $\subseteq \langle S\rangle$：由吸收性每项
$x_i s_i y_i\in\langle S\rangle$，有限和亦然。右端
$\supseteq \langle S\rangle$：右端对加减法封闭，且被任意环元素左乘
（右乘）后形状不变（乘进去的元素吸收进 $x_i$ 或 $y_i$），故右端本身是
包含 $S$ 的理想。
\end{proof}

\begin{definition}[理想的和]
设 $I, J$ 为理想，定义
\[
  I + J = \{\,x + y \mid x\in I,\ y\in J\,\}.
\]
\end{definition}

\begin{proposition}
$I + J$ 仍是理想。
\end{proposition}

\begin{proof}
加减封闭：$(x_1 + y_1)\pm(x_2 + y_2) = (x_1\pm x_2) + (y_1\pm y_2)\in I+J$。
吸收性：对任意 $r\in R$，
\[
  r(x+y) = rx + ry, \qquad rx\in I,\ ry\in J,
\]
故 $r(x+y)\in I+J$；另一侧同理。
\end{proof}

\begin{definition}[理想的积]
字面乘积集合 $\{xy \mid x\in I, y\in J\}$ 一般不是理想，定义 $IJ$ 为它
\keyword{生成的理想}。
\end{definition}

\begin{example}[$\Z$ 的理想都是主理想]
设 $I$ 是 $\Z$ 的理想，令 $d = \gcd\{\,a \mid a\in I\,\}$。辗转相减始终
在 $I$ 内部进行，故 $d\in I$，从而 $(d)\subseteq I$；另一方面 $I$ 中
元素都是 $d$ 的倍数，即 $I\subseteq(d)$。所以 $I = (d)$：$\Z$ 的每个
理想都是主理想。由此
\[
  (a) + (b) = \big(\gcd(a,b)\big), \qquad
  (a)\cap(b) = \big(\mathrm{lcm}(a,b)\big), \qquad
  (a)(b) = (ab).
\]
这些好性质只对整数环成立，一般的环没有这么好的事情。
\end{example}

\subsection{同构定理（续）}

\begin{definition}[同构]
既是同态又是双射的映射称为\keyword{同构}；同构的两个对象在代数意义上
等价。
\end{definition}

\block{商同态}
对任意理想 $I$，自然满同态
\[
  \pi\colon R\longrightarrow R/I, \qquad x\mapsto x + I,
\]
的核恰为 $I$。

\begin{theorem}[第三同构定理]
设 $I\subseteq J$ 是 $R$ 的理想，则
\[
  (R/I)\,\big/\,(J/I) \;\cong\; R/J.
\]
\end{theorem}

\begin{proof}
构造 $\varphi\colon R/I\to R/J$，$x + I\mapsto x + J$（由
$I\subseteq J$ 知良定义）。它是满射，且
\[
  \ker\varphi = \{\,x + I : x\in J\,\} = J/I,
\]
由第一同构定理即得。
\end{proof}

\begin{theorem}[对应定理]
设 $\varphi\colon R\to R'$ 是满同态，则有双射
\[
  \{\,R\text{ 中包含 }\ker\varphi\text{ 的理想}\,\}
  \;\longleftrightarrow\;
  \{\,R'\text{ 的理想}\,\},
\]
正向取像 $I\mapsto\varphi(I)$，反向取原像
$J'\mapsto\varphi^{-1}(J')$。特别地（取 $\varphi$ 为商同态），$R/I$
的理想与 $R$ 中包含 $I$ 的理想一一对应。
\end{theorem}
